%% =====================================================================
%%  Capítulo 9 -- Retroalimentación de estado, observadores y acción integral
%%  Migrado al formato nuevo (noviembre de 2026) desde Handout7_Ctrl.tex.
%%  Cambios de fondo respecto a la versión anterior (ver MIGRACION.md):
%%    el ESTIMADO del observador se escribe \tilde{x} (macro \est), porque
%%    \hat{x} es la variable desviación en todo el curso; el error de
%%    estimación es \varepsilon (e queda para el error de seguimiento);
%%    "Ackerman" -> Ackermann; signo del polinomio del servo corregido;
%%    K del ejemplo diag(-1,-3) calculado; Ke = 18 314 (exacto).
%% =====================================================================

\capitulo[cap:re]{Retroalimentación de estado y observadores}
  {[L]a virginidad o, mejor, la inocencia se pierde de otro modo. Se pierde, como se sabe, cuando entran en juego las emociones.}
  {Alberto Fuguet}

\begin{objetivos}[partes/09-retroalimentacion-estado.html]
  \item Explicar con controlabilidad y observabilidad por qué una función de transferencia puede esconder modos, y decidir si un sistema es estabilizable y detectable.
  \item Diseñar una retroalimentación de estado $u=r-\Km x$ que ubique los polos donde piden las especificaciones, por comparación de coeficientes y con la fórmula de Ackermann.
  \item Diseñar un observador de Luenberger y unirlo con la retroalimentación de estado usando el principio de separación.
  \item Agregar acción integral para anular el error en estado estacionario.
  \item Verificar cada diseño en \matlab/\simulink, incluido el precio que se paga en esfuerzo de control.
\end{objetivos}

%% ---------------------------------------------------------------------
\section{Lo que la función de transferencia esconde}
\label{sec:re-kalman}
%% ---------------------------------------------------------------------

Hasta ahora hemos diseñado controladores en el dominio de Laplace: un PID, un compensador ajustado sobre la función de transferencia. Volvemos ahora a la representación de estado del capítulo~4,
\begin{equation}
  \dot x=\Am x+\Bm u,\qquad y=\Cm x+\Dm u,
  \label{eq:re-estado}
\end{equation}
con $x\in\mathbb{R}^n$ (estado), $u\in\mathbb{R}^m$ (entradas), $y\in\mathbb{R}^p$ (salidas), $\Am$ de $n\times n$, $\Bm$ de $n\times m$, $\Cm$ de $p\times n$ y $\Dm$ de $p\times m$. La idea del capítulo es sencilla de enunciar: si medimos \emph{todo} el estado, podemos escoger $u$ como una combinación de los estados y poner los polos del lazo cerrado donde queramos. Antes de hacerlo hay que contestar dos preguntas: ¿la entrada alcanza a mover todos los estados? y ¿la salida alcanza a mostrarlos todos? Son la \clave{controlabilidad} y la \clave{observabilidad}.

\begin{historia}[Kalman y la teoría de sistemas de estado]
Rudolf E.\ Kalman introdujo los conceptos de controlabilidad y observabilidad, y su dualidad, en el primer congreso de la IFAC, en Moscú, en 1960~\citep{Kal60}. Una de sus motivaciones era explicar por qué fracasaba un método popular para controlar plantas inestables: cancelar los polos inestables con ceros del compensador. La cancelación algebraica podía ser perfecta y aun así el sistema real resultaba inestable. En 1963 mostró que todo sistema lineal se puede descomponer en cuatro partes según sea o no controlable y observable, y que la función de transferencia solo ``ve'' la parte que es ambas cosas~\citep{Kal63}.
\end{historia}

\begin{ejemplo}[label=ej:re-friedland]{Cuatro estados, una función de primer orden}
Considere el sistema~\citep{Fri86}
\[
\begin{aligned}
  \dot x_1&=2x_1+3x_2+2x_3+x_4+u, & \dot x_2&=-2x_1-3x_2-2u,\\
  \dot x_3&=-2x_1-2x_2-4x_3+2u,   & \dot x_4&=-2x_1-2x_2-2x_3-5x_4-u,
\end{aligned}
\qquad y=7x_1+6x_2+4x_3+2x_4.
\]
Calcule su función de transferencia.
\solucion $H(s)=\Cm(s\Id-\Am)^{-1}\Bm+\Dm$ da
\[
  H(s)=\frac{s^3+9s^2+26s+24}{s^4+10s^3+35s^2+50s+24}=\frac{(s+2)(s+3)(s+4)}{(s+1)(s+2)(s+3)(s+4)}=\frac{1}{s+1}.
\]
¡Un sistema de cuarto orden se ve desde afuera como uno de primer orden! Con un cambio de variables (una \emph{transformación de similaridad} que lleva $\Am$ a su forma diagonal) el mismo sistema queda
\[
  \dot z_1=-z_1+u,\qquad \dot z_2=-2z_2,\qquad \dot z_3=-3z_3+u,\qquad \dot z_4=-4z_4,\qquad y=z_1+z_2,
\]
y la \cref{fig:re-friedland} lo dice todo: $z_1$ es afectado por $u$ y visible en $y$; $z_2$ es visible pero $u$ no lo mueve; $z_3$ lo mueve $u$ pero no se ve en $y$; $z_4$ ni lo uno ni lo otro. Solo $z_1$ llega a la función de transferencia.
\end{ejemplo}

\begin{figure}[ht]
\centering
\includegraphics[width=0.85\linewidth]{cap09/modal_friedland}
\caption{Ejemplo~\ref{ej:re-friedland} en coordenadas modales. Solo el modo $z_1$ conecta la entrada con la salida; los otros tres existen, pero la función de transferencia no los muestra.}
\label{fig:re-friedland}
\end{figure}

\codigorepo{cap09_retroalimentacion_estado/C9_Ej1_Kalman.m}{kalman}{ejemplo \ref{ej:re-friedland}}

Kalman mostró que esto es general: todo sistema \eqref{eq:re-estado} se puede descomponer en (1) una parte controlable y observable, (2) una observable pero no controlable, (3) una controlable pero no observable y (4) una que no es ni lo uno ni lo otro~\citep{Kal63}. La función de transferencia está determinada solo por la primera. Un sistema que contiene una parte no controlable se dice \emph{no controlable}; uno que contiene una parte no observable, \emph{no observable}.

\begin{definicion}[label=def:re-estabilizable]{Estabilizable y detectable}
Un sistema es \clave{estabilizable} si todos sus modos no controlables son estables, y \clave{detectable} si todos sus modos no observables son estables.
\end{definicion}

En el ejemplo, los modos ``escondidos'' están en $-2$, $-3$ y $-4$: son estables, el sistema es estabilizable y detectable, y podemos ignorar esas dinámicas internas sin riesgo. Pero si uno de ellos estuviera en el semiplano derecho, ninguna medición de $y$ lo delataría y ningún diseño sobre $u$ lo corregiría: el sistema se desestabilizaría ``por dentro''.

\begin{cuidado}[Cancelar un polo inestable con un cero]
Cancelar un polo en el semiplano derecho con un cero del controlador produce, en el papel, una función de transferencia estable. En la realidad, el modo inestable sigue ahí: se volvió no controlable o no observable, y cualquier perturbación o condición inicial lo excita. Es exactamente lo que Kalman explicó. Nunca se cancelan polos inestables.
\end{cuidado}

%% ---------------------------------------------------------------------
\section{Retroalimentación de estado}
\label{sec:re-retro}
%% ---------------------------------------------------------------------

Supongamos que medimos todos los estados. La ley de control más sencilla que los usa es lineal:
\begin{equation}
  u(t)=r(t)-\Km\,x(t),
  \label{eq:re-ley}
\end{equation}
donde $r$ es una nueva entrada de referencia y $\Km$ (de $m\times n$; un vector fila en sistemas de una entrada) es lo \emph{único} que diseñamos (\cref{fig:re-lazo}). Reemplazando en \eqref{eq:re-estado},
\begin{equation}
  \dot x=(\Am-\Bm\Km)\,x+\Bm r,\qquad y=\Cm x.
  \label{eq:re-lazocerrado}
\end{equation}
El lazo cerrado es otro sistema de estado con matriz $\Am-\Bm\Km$: sus polos son las raíces de
\begin{equation}
  \det\bigl(\lambda\Id-(\Am-\Bm\Km)\bigr)=0,
  \label{eq:re-ecaracteristica}
\end{equation}
y escoger $\Km$ es escoger esa ecuación característica. Con $r=0$ el lazo es un \clave{regulador}: lleva el estado al equilibrio en el origen, y basta que $\Am-\Bm\Km$ sea Hurwitz (todos sus valores propios con parte real negativa). Con $r\neq0$ acotada, el punto de equilibrio se mueve, pero sigue siendo asintóticamente estable.

\begin{figure}[ht]
\centering
\includegraphics[width=0.8\linewidth]{cap09/lazo_estado}
\caption{Retroalimentación de estado $u=r-\Km x$. Con $r=0$ es un regulador.}
\label{fig:re-lazo}
\end{figure}

\begin{ejemplo}[label=ej:re-motivador]{Estabilizar una planta inestable}
Sea $\dot x=\begin{bmatrix}-1&a\\3&-2\end{bmatrix}x+\begin{bmatrix}0\\1\end{bmatrix}u$, $y=[1\;\;0]\,x$. Diseñe $u=-\Km x$ para que los polos queden en $-2$ y $-3$.
\solucion Con $\Km=[K_1\;\;K_2]$,
\[
  \Am-\Bm\Km=\begin{bmatrix}-1&a\\3-K_1&-2-K_2\end{bmatrix},\qquad
  \det\bigl(\lambda\Id-(\Am-\Bm\Km)\bigr)=\lambda^2+(3+K_2)\lambda+2+K_2+a(K_1-3).
\]
El polinomio deseado es $(\lambda+2)(\lambda+3)=\lambda^2+5\lambda+6$. Igualando coeficientes, $K_2=2$ y $K_1=3+\dfrac2a$ (con $a\neq0$).

Con $a=2$ la planta tiene valores propios en $1$ y $-4$: el equilibrio es inestable. Con $\Km=[4\;\;2]$ los estados regresan al origen (\cref{fig:re-motivador}); el precio es un esfuerzo inicial $u(0)=-\Km x(0)=-6$.
\end{ejemplo}

\begin{figure}[ht]
\centering
\includegraphics[width=\linewidth]{cap09/f_motivador}
\caption{Ejemplo~\ref{ej:re-motivador} con $a=2$ y $x(0)=[1\;\;1]^\top$: sin control los estados divergen; con $u=-\Km x$ vuelven al origen.}
\label{fig:re-motivador}
\end{figure}

\codigorepo{cap09_retroalimentacion_estado/C9_Ej2_Motivador.m}{motivador}{ejemplo \ref{ej:re-motivador}}

\begin{pausa}
En el ejemplo~\ref{ej:re-motivador}, ¿qué pasa si $a=0$? ¿Qué polo no se puede mover?
\respuesta{Con $a=0$ el polinomio queda $\lambda^2+(3+K_2)\lambda+2+K_2=(\lambda+1)(\lambda+2+K_2)$: el factor $(\lambda+1)$ no depende de $\Km$. El polo en $-1$ no se mueve, porque con $a=0$ la ecuación de $x_1$ es $\dot x_1=-x_1$ y la entrada no llega a $x_1$. Es el tema de la próxima sección.}
\end{pausa}

%% ---------------------------------------------------------------------
\section{Controlabilidad}
\label{sec:re-controlabilidad}
%% ---------------------------------------------------------------------

La pausa anterior muestra que no siempre se pueden poner los polos donde uno quiera. ¿Cuándo sí?

\begin{definicion}[label=def:re-controlable]{Controlabilidad}
El sistema $\dot x=\Am x+\Bm u$ es \clave{controlable} si, para \emph{cualquier} estado inicial $x(0)=x_0$ y \emph{cualquier} estado final $x_T$, existe una entrada $u(t)$ que lleva el sistema de $x_0$ a $x_T$ en un tiempo \emph{finito} $T$.
\end{definicion}

Hay una caracterización equivalente que es la que nos interesa para diseñar: el sistema es controlable si y solo si, para cualquier polinomio mónico $\alpha_c(\lambda)$ de grado $n$ con coeficientes reales, existe $\Km$ tal que el polinomio característico de $\Am-\Bm\Km$ es $\alpha_c(\lambda)$ (\emph{ubicación arbitraria de polos}). En sistemas de una entrada ese $\Km$ es, además, único. Y hay una prueba algebraica que no exige calcular nada más que productos de matrices.

\begin{definicion}[label=def:re-ctrb]{Matriz de controlabilidad}
El sistema con $\Am\in\mathbb{R}^{n\times n}$ es controlable si y solo si la \clave{matriz de controlabilidad}
\begin{equation}
  \Ctrb=\bigl[\,\Bm\;\;\Am\Bm\;\;\Am^2\Bm\;\;\cdots\;\;\Am^{n-1}\Bm\,\bigr]
  \label{eq:re-ctrb}
\end{equation}
tiene rango $n$. Con una sola entrada $\Ctrb$ es cuadrada y basta revisar que $\det\Ctrb\neq0$.
\end{definicion}

La intuición: $\Bm u$ empuja el estado en la dirección de $\Bm$; un instante después la dinámica ha llevado ese empujón a la dirección $\Am\Bm$, luego a $\Am^2\Bm$, y así. Si esas $n$ direcciones generan todo $\mathbb{R}^n$, la entrada puede llevar el estado a cualquier parte. En el ejemplo~\ref{ej:re-motivador}, $\Ctrb=\begin{bmatrix}0&a\\1&-2\end{bmatrix}$ y $\det\Ctrb=-a$: controlable si y solo si $a\neq0$.

\begin{ejemplo}[label=ej:re-diag]{Dos modos y una entrada}
Para $\Am=\begin{bmatrix}-1&0\\0&-3\end{bmatrix}$, $\Bm=\begin{bmatrix}1\\1\end{bmatrix}$: ¿es controlable? Si lo es, halle $\Km$ para tener los dos polos en $-2$.
\solucion $\Ctrb=[\Bm\;\;\Am\Bm]=\begin{bmatrix}1&-1\\1&-3\end{bmatrix}$, con $\det\Ctrb=-2\neq0$: controlable (la entrada llega a los dos modos). Los polos de la planta son $-1$ y $-3$. Con $\Km=[K_1\;\;K_2]$,
\[
  \det\bigl(\lambda\Id-(\Am-\Bm\Km)\bigr)=\lambda^2+(4+K_1+K_2)\lambda+3+3K_1+K_2,
\]
e igualando con $(\lambda+2)^2=\lambda^2+4\lambda+4$: $K_1+K_2=0$ y $3K_1+K_2=1$, de donde $\Km=[0.5\;\;-0.5]$.
\end{ejemplo}

\codigorepo{cap09_retroalimentacion_estado/C9_Ej3_Controlabilidad.m}{rango}{ejemplo \ref{ej:re-diag}}

\begin{profundiza}[¿Cuál modo falla? La prueba PBH]
El rango de $\Ctrb$ dice \emph{si} el sistema es controlable, pero no \emph{qué} modo falla. La prueba de Popov-Belevitch-Hautus lo dice: $(\Am,\Bm)$ es controlable si y solo si
\[
  \operatorname{rango}\bigl[\,\Am-\lambda\Id\;\;\Bm\,\bigr]=n\quad\text{para cada valor propio }\lambda\text{ de }\Am.
\]
El valor propio para el cual el rango cae es el modo no controlable. En el ejemplo~\ref{ej:re-motivador} con $a=0$, el rango cae en $\lambda=-1$. Y el sistema es estabilizable si y solo si los $\lambda$ que fallan tienen parte real negativa.
\end{profundiza}
\codigorepo{cap09_retroalimentacion_estado/C9_Ej3_Controlabilidad.m}{pbh}{prueba PBH}

\subsection{Controlable no quiere decir barato}

La controlabilidad es una propiedad de sí o no, pero en la práctica importa \emph{cuánto}. En el ejemplo~\ref{ej:re-motivador}, $K_1=3+2/a$: cuando $a\to0$ el sistema sigue siendo controlable, pero la ganancia necesaria crece sin límite, y con ella el esfuerzo de control (\cref{fig:re-condicion}). El número de condición de $\Ctrb$ es un buen termómetro: si es muy grande, $\Ctrb$ está ``casi'' sin rango completo y mover los polos sale caro (actuadores que se saturan, ruido amplificado).

\begin{figure}[ht]
\centering
\includegraphics[width=0.85\linewidth]{cap09/f_condicion}
\caption{Ejemplo~\ref{ej:re-motivador} con los polos siempre en $-2$ y $-3$: a medida que $a\to0$, el número de condición de $\Ctrb$ y el pico de $|u|$ crecen como $1/a$.}
\label{fig:re-condicion}
\end{figure}

\codigorepo{cap09_retroalimentacion_estado/C9_Ej3_Controlabilidad.m}{condicion}{\cref{fig:re-condicion}}

%% ---------------------------------------------------------------------
\section{La fórmula de Ackermann}
\label{sec:re-ackermann}
%% ---------------------------------------------------------------------

Comparar coeficientes funciona bien con $n=2$ o $3$, pero con más estados el álgebra se vuelve pesada. Para sistemas de una entrada hay una fórmula cerrada.

\begin{historia}[Jürgen Ackermann]
La fórmula lleva el nombre del ingeniero alemán Jürgen Ackermann, que la publicó en 1972 en un artículo sobre el diseño de sistemas de control lineales en el espacio de estado~\citep{Ack72}. Ackermann trabajó luego en el Centro Aeroespacial Alemán (DLR) en control robusto, con aplicaciones a la dirección de automóviles.
\end{historia}

\begin{definicion}[label=def:re-ackermann]{Fórmula de Ackermann}
Si el sistema de una entrada es controlable y el polinomio característico deseado es $q(\lambda)=\lambda^n+\alpha_{n-1}\lambda^{n-1}+\cdots+\alpha_1\lambda+\alpha_0$, entonces
\begin{equation}
  \Km=\bigl[\,0\;\;0\;\;\cdots\;\;0\;\;1\,\bigr]\,\Ctrb^{-1}\,q(\Am),
  \qquad
  q(\Am)=\Am^n+\alpha_{n-1}\Am^{n-1}+\cdots+\alpha_1\Am+\alpha_0\Id.
  \label{eq:re-ackermann}
\end{equation}
\end{definicion}

Observe que $q(\Am)$ es el polinomio deseado evaluado en la \emph{matriz} $\Am$ (con $\Id$ en lugar del término independiente), y que la fórmula exige $\Ctrb^{-1}$: solo funciona si el sistema es controlable.

\begin{ejemplo}[label=ej:re-ackermann]{El doble integrador}
Para $Y(s)/U(s)=1/s^2$, halle $\Km$ para tener los polos en $-1\pm j$.
\solucion En variables de fase, $\Am=\begin{bmatrix}0&1\\0&0\end{bmatrix}$, $\Bm=\begin{bmatrix}0\\1\end{bmatrix}$. El polinomio deseado es $q(\lambda)=\lambda^2+2\lambda+2$. Entonces
\[
  \Ctrb=[\Bm\;\;\Am\Bm]=\begin{bmatrix}0&1\\1&0\end{bmatrix}=\Ctrb^{-1},\qquad
  q(\Am)=\Am^2+2\Am+2\Id=\begin{bmatrix}2&2\\0&2\end{bmatrix},
\]
\[
  \Km=[0\;\;1]\begin{bmatrix}0&1\\1&0\end{bmatrix}\begin{bmatrix}2&2\\0&2\end{bmatrix}=[1\;\;0]\begin{bmatrix}2&2\\0&2\end{bmatrix}=[2\;\;2].
\]
Comprobación: $\Am-\Bm\Km=\begin{bmatrix}0&1\\-2&-2\end{bmatrix}$, cuyo polinomio es $\lambda^2+2\lambda+2$.
\end{ejemplo}

\codigorepo{cap09_retroalimentacion_estado/C9_Ej4_Ackermann.m}{ackermann}{ejemplo \ref{ej:re-ackermann}: la fórmula paso a paso, \cmd{acker} y \cmd{place}}

\begin{cuidado}[\texttt{acker} o \texttt{place}]
\cmd{acker} implementa \eqref{eq:re-ackermann}: es exacta en papel pero numéricamente frágil para $n$ grande (invierte $\Ctrb$, que puede estar mal condicionada). \cmd{place} usa un algoritmo más robusto, sirve para varias entradas, pero no acepta polos repetidos más veces que el número de entradas. Para los ejemplos del curso dan lo mismo.
\end{cuidado}

\begin{receta}[Ubicación de polos por retroalimentación de estado]
\begin{enumerate}
  \item Escribir el modelo de estado y verificar que $\operatorname{rango}\Ctrb=n$.
  \item Traducir las especificaciones en polos deseados (\cref{sec:re-especificaciones}) y armar $q(\lambda)$.
  \item Calcular $\Km$: comparando los coeficientes de $\det(\lambda\Id-\Am+\Bm\Km)$ con $q(\lambda)$, o con \eqref{eq:re-ackermann}.
  \item Verificar \cmd{eig(A-B*K)}, simular, y mirar el esfuerzo de control $u(t)$.
\end{enumerate}
\end{receta}

\begin{pausa}
Use la fórmula de Ackermann para $1/(s(s+3))$ en variables de fase, con polos en $-4\pm j4$.
\respuesta{$\Am=\begin{bmatrix}0&1\\0&-3\end{bmatrix}$, $\Ctrb=\begin{bmatrix}0&1\\1&-3\end{bmatrix}$, $\Ctrb^{-1}=\begin{bmatrix}3&1\\1&0\end{bmatrix}$, $q(\Am)=\Am^2+8\Am+32\Id=\begin{bmatrix}32&5\\0&17\end{bmatrix}$. Entonces $\Km=[1\;\;0]\,q(\Am)=[32\;\;5]$. Comprobación: $\lambda^2+(3+K_2)\lambda+K_1=\lambda^2+8\lambda+32$.}
\end{pausa}

%% ---------------------------------------------------------------------
\section{De las especificaciones a los polos}
\label{sec:re-especificaciones}
%% ---------------------------------------------------------------------

Con $n$ estados hay que escoger $n$ polos. La estrategia es la del \cref{cap:so}: dos polos \emph{dominantes} fijados por las especificaciones del segundo orden ($\OS$ fija $\zeta$; $\tp$ o $\ts$ fijan $\wn$) y los demás bastante más a la izquierda, para que sus modos mueran rápido. Una regla usual es ponerlos entre cinco y diez veces más lejos que la parte real de los dominantes. Si la planta tiene un cero de fase mínima cerca de los dominantes, conviene usar uno de los polos sobrantes para \emph{cancelarlo} (sin modificar la ganancia).

\begin{ejemplo}[label=ej:re-trespolos]{Un tercer orden con especificaciones}
La planta $G_p(s)=\dfrac{10}{(s+10)(s^2+0.5s+1)}$ tiene $\OS=44\%$ y $\tp=3.35$\,s. Diseñe $u=r-\Km x$ para $\OS=10\%$ y $\tp=1$\,s.
\solucion $\OS=10\%$ da $\zeta=0.591$, y $\tp=1$ da $\wn=\pi/\sqrt{1-\zeta^2}=3.90$\,rad/s: polos dominantes en $-2.30\pm j3.14$. El tercero va diez veces más lejos, en $-23.0$. Con la representación que entrega \cmd{tf2ss},
\[
  \Km=[17.131\;\;115.21\;\;339.34].
\]
La respuesta (\cref{fig:re-trespolos}) tiene $\OS=9.8\%$ y $\tp=1.05$\,s: el tercer polo, aunque lejano, la desvía un poco de las fórmulas. Pero el valor final es $0.029$: la salida no sigue la referencia. Lo arreglamos en la \cref{sec:re-integral}.
\end{ejemplo}

\begin{figure}[ht]
\centering
\includegraphics[width=\linewidth]{cap09/f_tres_polos}
\caption{Ejemplo~\ref{ej:re-trespolos}. La forma de la respuesta cumple; el valor final (que a la derecha se normalizó para comparar la forma), no.}
\label{fig:re-trespolos}
\end{figure}

\codigorepo{cap09_retroalimentacion_estado/C9_Ej5_Especificaciones.m}{tercerorden}{ejemplo \ref{ej:re-trespolos}}

\begin{cuidado}[Las ganancias dependen de la representación; los polos, no]
El mismo sistema tiene infinitas representaciones de estado, y $\Km$ cambia con cada una porque los estados son otros. \cmd{tf2ss} entrega la forma con la primera fila $[-a_{n-1}\;\cdots\;-a_0]$; las \emph{variables de fase} del capítulo de espacio de estado tienen la última fila $[-a_0\;\cdots\;-a_{n-1}]$. En el ejemplo~\ref{ej:re-trespolos}, en variables de fase sale $\Km=[339.34\;\;115.21\;\;17.131]$: las mismas cifras al revés, porque los estados están numerados al revés. Los polos del lazo cerrado son idénticos. Diga siempre en qué representación trabaja; la función del curso \cmd{variablesFase} construye la de las notas.
\end{cuidado}

\subsection{Cuando la referencia entra por otro lado}

No siempre la referencia se suma a $u$. Si la salida es un estado (por ejemplo $y=x_1$), lo natural es compararla con $r$ como en un lazo clásico.

\begin{ejemplo}[label=ej:re-servo]{Un servo de tercer orden}
La planta es $\dfrac{Y(s)}{U(s)}=\dfrac{1}{s(s+1)(s+2)}$ con la configuración de la \cref{fig:re-servo}: $u=K_1(r-x_1)-K_2x_2-K_3x_3$. Ubique los polos en $-2\pm j3.464$ y $-10$.
\solucion En variables de fase, $\Am=\begin{bmatrix}0&1&0\\0&0&1\\0&-2&-3\end{bmatrix}$, $\Bm=\begin{bmatrix}0\\0\\1\end{bmatrix}$, $\Cm=[1\;\;0\;\;0]$. La matriz de controlabilidad,
\[
  \Ctrb=[\Bm\;\;\Am\Bm\;\;\Am^2\Bm]=\begin{bmatrix}0&0&1\\0&1&-3\\1&-3&7\end{bmatrix},
\]
tiene rango $3$. La ley es $u=-\Km x+K_1r$ con $\Km=[K_1\;\;K_2\;\;K_3]$, y el polinomio característico es
\[
  \lambda^3+(3+K_3)\lambda^2+(2+K_2)\lambda+K_1.
\]
El deseado es $(\lambda+10)(\lambda^2+4\lambda+16)=\lambda^3+14\lambda^2+56\lambda+160$, luego $\Km=[160\;\;54\;\;11]$ y
\[
  \dot x=\begin{bmatrix}0&1&0\\0&0&1\\-160&-56&-14\end{bmatrix}x+\begin{bmatrix}0\\0\\160\end{bmatrix}r,\qquad y=[1\;\;0\;\;0]\,x.
\]
Ahora el valor final \emph{sí} es $1$: la planta tiene un integrador y la referencia se compara con la salida, como en un lazo tipo 1 (\cref{cap:er}).
\end{ejemplo}

\begin{figure}[ht]
\centering
\includegraphics[width=0.9\linewidth]{cap09/servo_k1}
\caption{Ejemplo~\ref{ej:re-servo}: la referencia se compara con $x_1=y$; los otros estados se retroalimentan con $K_2$ y $K_3$.}
\label{fig:re-servo}
\end{figure}

\codigorepo{cap09_retroalimentacion_estado/C9_Ej6_Servo.m}{servo}{ejemplo \ref{ej:re-servo}}

\begin{pausa}
Para $T(s)=\dfrac{100(s+1)}{s(s+2)(s+3)}$ se pide $\OS=10\%$ y $\tp=1$\,s con $u=r-\Km x$, cancelando el cero con el tercer polo. ¿Dónde va el tercer polo y qué función de transferencia queda en lazo cerrado?
\respuesta{En $-1$, sobre el cero. Queda $\dfrac{100}{s^2+4.605s+15.17}$: la forma pedida ($\zeta=0.591$, $\wn=3.90$), pero con valor final $100/15.17=6.59$, no $1$. Otra vez hace falta la acción integral.}
\end{pausa}

%% ---------------------------------------------------------------------
\section{Observadores de estado}
\label{sec:re-observador}
%% ---------------------------------------------------------------------

Todo lo anterior supone que medimos todos los estados. Casi nunca es así: medir cuesta, y algunos estados (una velocidad interna, una concentración) no tienen sensor. La salida de este problema es \emph{estimarlos}: construir dentro del controlador una copia del modelo que, alimentada con la entrada y con la salida medida, reproduzca el estado.

\begin{historia}[David Luenberger]
David G.\ Luenberger acuñó el término \emph{observador} en su tesis doctoral en Stanford (1963) y lo publicó en 1964~\citep{Lue64}. Mostró que para cualquier sistema lineal observable se puede construir un observador cuyo error de estimación tienda a cero tan rápido como se quiera, y que su diseño es el ``dual'' de la ubicación de polos por retroalimentación de estado.
\end{historia}

\paragraph{La idea}
Una copia del modelo, $\dot{\est{x}}=\Am\est{x}+\Bm u$, reproduciría el estado si arrancara exactamente en $x(0)$ y el modelo fuera perfecto. Como no es así, la corregimos con lo que sí sabemos: la diferencia entre la salida medida $y$ y la que predice la copia, $\Cm\est{x}$ (\cref{fig:re-observador}):
\begin{equation}
  \dot{\est{x}}=\Am\est{x}+\Bm u+\Lm\,(y-\Cm\est{x})=(\Am-\Lm\Cm)\,\est{x}+\Bm u+\Lm y.
  \label{eq:re-observador}
\end{equation}
$\Lm$ (de $n\times p$) es la \clave{ganancia del observador}. Restando \eqref{eq:re-observador} de \eqref{eq:re-estado}, el \clave{error de estimación} $\errest=x-\est{x}$ cumple
\begin{equation}
  \dot{\errest}=(\Am-\Lm\Cm)\,\errest.
  \label{eq:re-errorobs}
\end{equation}
Si $\Am-\Lm\Cm$ es Hurwitz, $\est{x}\to x$ \emph{sin importar la entrada}, y la velocidad la fijan los valores propios de $\Am-\Lm\Cm$, que escogemos con $\Lm$.

\begin{figure}[ht]
\centering
\includegraphics[width=0.8\linewidth]{cap09/observador}
\caption{Observador de Luenberger: una copia del modelo corregida con el error de salida $y-\Cm\est{x}$.}
\label{fig:re-observador}
\end{figure}

\begin{cuidado}[$\est{x}$ y no $\hat{x}$]
En muchos libros el estimado se escribe $\hat{x}$. En este curso $\hat{x}$ es la \emph{variable desviación} (\cref{cap:so}), y los dos símbolos aparecen juntos cuando se diseña un observador para un modelo linealizado. Por eso el estimado es $\est{x}$ y el error de estimación es $\errest$ (la letra $e$ queda para el error de seguimiento $r-y$).
\end{cuidado}

\begin{definicion}[label=def:re-observable]{Observabilidad}
El sistema \eqref{eq:re-estado} es \clave{observable} si el estado inicial $x(0)$ (y con él todo $x(t)$) se puede determinar a partir de la entrada $u(\tau)$ y la salida $y(\tau)$ en un intervalo finito $0\leq\tau\leq t$. Equivalentemente, si y solo si la \clave{matriz de observabilidad}
\begin{equation}
  \Obsv=\begin{bmatrix}\Cm\\ \Cm\Am\\ \vdots\\ \Cm\Am^{n-1}\end{bmatrix}
  \label{eq:re-obsv}
\end{equation}
tiene rango $n$. Y los polos del observador (los valores propios de $\Am-\Lm\Cm$) se pueden ubicar arbitrariamente si y solo si $(\Am,\Cm)$ es observable.
\end{definicion}

\paragraph{Dualidad}
Los valores propios de $\Am-\Lm\Cm$ son los de su transpuesta, $\Am^\top-\Cm^\top\Lm^\top$, que tiene la forma $\Am-\Bm\Km$ con $\Am\to\Am^\top$, $\Bm\to\Cm^\top$ y $\Km\to\Lm^\top$. Diseñar un observador es ubicar polos en el sistema ``dual''. En \matlab: \cmd{L = acker(A',C',p)'} (o \cmd{place}). La fórmula de Ackermann para el observador es
\begin{equation}
  \Lm=\rho(\Am)\,\Obsv^{-1}\begin{bmatrix}0\\ \vdots\\ 0\\ 1\end{bmatrix},\qquad
  \rho(\Am)=\Am^n+\beta_{n-1}\Am^{n-1}+\cdots+\beta_0\Id,
  \label{eq:re-ackermannobs}
\end{equation}
donde $\rho(\lambda)=\lambda^n+\beta_{n-1}\lambda^{n-1}+\cdots+\beta_0$ es el polinomio deseado del observador.

\begin{ejemplo}[label=ej:re-observador]{Un observador con Ackermann}
Para $\Am=\begin{bmatrix}2&3\\-1&4\end{bmatrix}$, $\Cm=[1\;\;0]$, halle $\Lm$ para que el observador tenga $\rho(\lambda)=\lambda^2+2\zeta\wn\lambda+\wn^2$ con $\zeta=0.8$ y $\wn=10$.
\solucion $\rho(\lambda)=\lambda^2+16\lambda+100$. Entonces
\[
  \rho(\Am)=\Am^2+16\Am+100\Id=\begin{bmatrix}133&66\\-22&177\end{bmatrix},\qquad
  \Obsv=\begin{bmatrix}1&0\\2&3\end{bmatrix},\qquad
  \Obsv^{-1}=\begin{bmatrix}1&0\\-2/3&1/3\end{bmatrix},
\]
y $\Lm=\rho(\Am)\Obsv^{-1}[0\;\;1]^\top=\rho(\Am)\,[0\;\;\tfrac13]^\top=[22\;\;59]^\top$. Con la planta $\Am=\begin{bmatrix}0&1\\-10&-2\end{bmatrix}$ y el mismo $\rho$, $\Obsv=\Id$ y sale $\Lm=[14\;\;62]^\top$; la \cref{fig:re-obs} muestra cómo el estimado alcanza al estado real partiendo de condiciones iniciales distintas.
\end{ejemplo}

\begin{figure}[ht]
\centering
\includegraphics[width=\linewidth]{cap09/f_observador}
\caption{Ejemplo~\ref{ej:re-observador} con $\Am=\begin{bmatrix}0&1\\-10&-2\end{bmatrix}$: la planta arranca en $x(0)=[1\;\;0]^\top$ y el observador en cero. En medio segundo el error de estimación ya es menor que el $4\%$.}
\label{fig:re-obs}
\end{figure}

\codigorepo{cap09_retroalimentacion_estado/C9_Ej7_Observador.m}{observador}{ejemplo \ref{ej:re-observador}}

\begin{ideaclave}
¿Qué tan rápido debe ser el observador? Lo bastante para que el estimado esté listo antes de que el controlador lo necesite: es usual ponerlo de dos a cinco veces más rápido que los polos de control. Más rápido no es gratis: un $\Lm$ grande amplifica el ruido de medición que entra por $y$. Es el mismo compromiso entre velocidad y ruido de todo el curso.
\end{ideaclave}

%% ---------------------------------------------------------------------
\section{Retroalimentación de estado con observador}
\label{sec:re-separacion}
%% ---------------------------------------------------------------------

El diseño completo se hace en cuatro pasos: (1) diseñar $\Km$ como si todos los estados se midieran; (2) diseñar un observador para los que no se miden; (3) conectar los dos, usando el estimado en lugar del estado:
\begin{equation}
  u=-\Km\est{x}\;(+Nr),
  \label{eq:re-leyobs}
\end{equation}
y (4) agregar la entrada de referencia. Al conjunto ``observador más $\Km$'' se le llama \clave{compensador} (\cref{fig:re-compensador}).

\begin{figure}[ht]
\centering
\includegraphics[width=0.9\linewidth]{cap09/compensador}
\caption{Compensador por variables de estado: el observador produce $\est{x}$ y la ley de control lo usa en lugar de $x$.}
\label{fig:re-compensador}
\end{figure}

¿Sigue siendo estable el lazo si usamos $\est{x}$ en lugar de $x$? Con $\est{x}=x-\errest$, la ley \eqref{eq:re-leyobs} (con $r=0$) da $\dot x=(\Am-\Bm\Km)x+\Bm\Km\errest$, y el error sigue cumpliendo \eqref{eq:re-errorobs}. Juntando,
\begin{equation}
  \begin{bmatrix}\dot x\\ \dot{\errest}\end{bmatrix}=
  \begin{bmatrix}\Am-\Bm\Km & \Bm\Km\\ 0 & \Am-\Lm\Cm\end{bmatrix}
  \begin{bmatrix}x\\ \errest\end{bmatrix}.
  \label{eq:re-aumentado}
\end{equation}
La matriz es triangular por bloques, así que su ecuación característica es el producto
\begin{equation}
  \det\bigl(\lambda\Id-(\Am-\Bm\Km)\bigr)\cdot\det\bigl(\lambda\Id-(\Am-\Lm\Cm)\bigr)=0.
\end{equation}

\begin{ideaclave}
\textbf{Principio de separación.} Los polos del lazo con observador son la unión de los polos de control (los de $\Am-\Bm\Km$) y los del observador (los de $\Am-\Lm\Cm$). Por eso $\Km$ y $\Lm$ se diseñan \emph{por separado}: si cada parte es estable, el conjunto lo es.
\end{ideaclave}

Visto desde afuera, el compensador es simplemente una función de transferencia entre la salida medida y la entrada a la planta (con $r=0$):
\begin{equation}
  \frac{U(s)}{Y(s)}=-\Km\bigl(s\Id-(\Am-\Bm\Km-\Lm\Cm)\bigr)^{-1}\Lm.
  \label{eq:re-compensador}
\end{equation}

\begin{ejemplo}[label=ej:re-separacion]{Un compensador completo}
Para $\Am=\begin{bmatrix}0&1\\-2&-3\end{bmatrix}$, $\Bm=\begin{bmatrix}0\\1\end{bmatrix}$, $\Cm=[1\;\;0]$ (la planta $1/((s+1)(s+2))$, de la que solo se mide la posición), diseñe $\Km$ para polos en $-2\pm j2$ y un observador con polos dobles en $-10$.
\solucion Control: $\lambda^2+(3+K_2)\lambda+2+K_1=\lambda^2+4\lambda+8$, luego $\Km=[6\;\;1]$. Observador: $\Am-\Lm\Cm=\begin{bmatrix}-L_1&1\\-2-L_2&-3\end{bmatrix}$ tiene polinomio $\lambda^2+(3+L_1)\lambda+3L_1+2+L_2$; igualando con $(\lambda+10)^2$, $\Lm=[17\;\;47]^\top$. Los cuatro polos de \eqref{eq:re-aumentado} son $-2\pm j2$, $-10$ y $-10$. Y el compensador \eqref{eq:re-compensador} resulta
\[
  \frac{U(s)}{Y(s)}=-\frac{149s+554}{s^2+21s+123},
\]
un compensador de segundo orden con un cero en $-3.72$ y polos en $-10.5\pm j3.57$: un ``adelanto con filtro'' que salió del espacio de estado sin que lo buscáramos.
\end{ejemplo}

\codigorepo{cap09_retroalimentacion_estado/C9_Ej8_Compensador.m}{separacion}{ejemplo \ref{ej:re-separacion}}

\begin{figure}[ht]
\centering
\includegraphics[width=0.85\linewidth]{cap09/f_separacion}
\caption{Ejemplo~\ref{ej:re-separacion} con la referencia escalada por $N$ para que el valor final sea $1$. Con un error inicial de estimación, el observador rápido (polos en $-10$) casi no se nota; uno lento (en $-4$) deforma el transitorio.}
\label{fig:re-separacionfig}
\end{figure}

\begin{modelo}[Observador.slx]
El script \cmd{crearModelo\_Observador.m} construye la planta (que arranca en $x(0)=[1\;\;-1]^\top$), el observador con entradas $[u;\,y]$ y la ley $u=-\Km\est{x}+Nr$. Un interruptor manual cambia $\est{x}$ por el estado real para comparar. Mueva los polos del observador (\cmd{po}) y observe cuánto tarda $\est{x}$ en alcanzar a $x$.
\end{modelo}

\begin{cuidado}[La separación no garantiza robustez]
El principio de separación supone que el modelo es exacto. Si la planta real difiere del modelo con que se diseñó el observador, el error de estimación ya no obedece \eqref{eq:re-errorobs} y los márgenes de estabilidad del compensador completo pueden ser mucho menores que los de la retroalimentación de estado sola. Hay que verificar el diseño con la planta perturbada (problema~\ref{prob:re-robustez}).
\end{cuidado}

%% ---------------------------------------------------------------------
\section{Acción integral: error nulo en estado estacionario}
\label{sec:re-integral}
%% ---------------------------------------------------------------------

Los ejemplos~\ref{ej:re-trespolos} y la pausa de la \cref{sec:re-especificaciones} mostraron el problema: $u=r-\Km x$ pone los polos donde queremos, pero nada obliga a que $y$ llegue a $r$. En el \cref{cap:er} vimos que el remedio es un integrador en el lazo (\cref{sec:er-estado}). En espacio de estado se hace igual.

\begin{ejemplo}[label=ej:re-sinintegral]{El error de la retroalimentación de estado pura}
Para $\Am=\begin{bmatrix}0&1\\-3&-5\end{bmatrix}$, $\Bm=\begin{bmatrix}0\\1\end{bmatrix}$, $\Cm=[1\;\;0]$, diseñe $u=r-\Km x$ con $\OS=10\%$ y $\ts=0.5$\,s, y calcule el error en estado estacionario ante un escalón unitario.
\solucion $\zeta=0.591$ y $\zeta\wn=4/0.5=8$, así que $\wn=13.53$ y el polinomio deseado es $s^2+16s+183.1$ (polos en $-8\pm j10.91$). El lazo cerrado tiene $s^2+(5+K_2)s+(3+K_1)$, luego $\Km=[180.1\;\;11]$. Con $\bar{\Am}=\Am-\Bm\Km$, $Y(s)/R(s)=\Cm(s\Id-\bar{\Am})^{-1}\Bm=1/(s^2+16s+183.1)$ y
\[
  \ess=\lim_{s\to0}s\,\frac1s\bigl(1-\Cm(s\Id-\bar{\Am})^{-1}\Bm\bigr)=1-\frac{1}{183.1}=0.9945.
\]
¡La salida llega a menos del $1\%$ de la referencia!
\end{ejemplo}

\paragraph{El estado integrador}
Agregamos un estado nuevo, la integral del error de seguimiento (\cref{fig:re-integral}):
\begin{equation}
  z(t)=\int_0^t e(\tau)\,\dif\tau,\qquad \dot z=e=r-\Cm x,
\end{equation}
y la ley de control $u=-\Km x+K_e z=-\bigl[\Km\;\;-K_e\bigr]\begin{bmatrix}x\\z\end{bmatrix}$. El sistema aumentado es
\begin{equation}
  \begin{bmatrix}\dot x\\ \dot z\end{bmatrix}=
  \begin{bmatrix}\Am-\Bm\Km & \Bm K_e\\ -\Cm & 0\end{bmatrix}
  \begin{bmatrix}x\\ z\end{bmatrix}+\begin{bmatrix}0\\1\end{bmatrix}r,
  \qquad y=\bigl[\Cm\;\;0\bigr]\begin{bmatrix}x\\ z\end{bmatrix}.
  \label{eq:re-integral}
\end{equation}
Si el lazo aumentado es estable, en estado estacionario $\dot z=0$, es decir $y=r$: el error es nulo para \emph{cualquier} valor de los parámetros de la planta, mientras el lazo siga estable. Hay un polo más que ubicar: si no hay ceros cerca, de cinco a diez veces más a la izquierda que los dominantes; si los hay, sobre el cero.

\begin{figure}[ht]
\centering
\includegraphics[width=0.85\linewidth]{cap09/integral_estado}
\caption{Retroalimentación de estado con acción integral: $u=-\Km x+K_ez$, con $\dot z=r-y$.}
\label{fig:re-integral}
\end{figure}

\begin{ejemplo}[label=ej:re-integral]{Ahora con integrador}
Repita el ejemplo~\ref{ej:re-sinintegral} con acción integral, con el tercer polo en $-100$.
\solucion El polinomio característico de \eqref{eq:re-integral} es $s^3+(5+K_2)s^2+(3+K_1)s+K_e$, y el deseado, $(s+100)(s^2+16s+183.1)=s^3+116s^2+1783.1s+18\,310$. Luego
\[
  K_1=1780.1,\qquad K_2=111,\qquad K_e=18\,314
\]
(con $\wn^2$ sin redondear, $183.137$; con $183.1$ saldría $18\,310$). La \cref{fig:re-integralfig} muestra que el error se anula. El precio es el esfuerzo: un tercer polo tan rápido pide picos de $u$ de más de cien veces el escalón.
\end{ejemplo}

\begin{figure}[ht]
\centering
\includegraphics[width=\linewidth]{cap09/f_integral}
\caption{Ejemplos~\ref{ej:re-sinintegral} y~\ref{ej:re-integral}: sin integrador la salida no sigue la referencia; con él la sigue, a cambio de un esfuerzo de control mucho mayor.}
\label{fig:re-integralfig}
\end{figure}

\codigorepo{cap09_retroalimentacion_estado/C9_Ej9_Integral.m}{integrador}{ejemplo \ref{ej:re-integral}}

\begin{cuidado}[Cuando la salida combina estados]
Si $\Cm$ tiene más de un elemento no nulo (por ejemplo $\Cm=[4\;\;1]$), $K_e$ aparece en \emph{varios} coeficientes del polinomio aumentado, no solo en el independiente. Calcule siempre $\det(\lambda\Id-\Am_{aum})$ completo; no copie la forma del ejemplo anterior (problema~\ref{prob:re-integralcero}).
\end{cuidado}

\begin{modelo}[RetroEstado.slx]
El script \cmd{crearModelo\_RetroEstado.m} diseña las dos leyes del ejemplo y construye dos copias de la planta: una con $u=r-\Km x$ y otra con $u=-\Km x+K_ez$. Los \emph{Scope} comparan las salidas y los esfuerzos. Pruebe a cambiar el tercer polo y a perturbar $\Am$ (la planta real): la versión con integrador sigue llegando a $1$.
\end{modelo}

\begin{profundiza}[Retroalimentación de estado con integrador = un PID]
Para una planta de segundo orden en variables de fase, $x_1=y$ y $x_2=\dot y$. La ley $u=-K_1y-K_2\dot y+K_e\int e$ es exactamente un controlador \textbf{I-PD}: la acción integral actúa sobre el error y las acciones proporcional y derivativa, solo sobre la salida (capítulo del PID). La formulación de estado agrega dos cosas: una manera sistemática de ubicar \emph{todos} los polos, y la posibilidad de usar un observador cuando $\dot y$ no se mide o la planta es de orden mayor (problema~\ref{prob:re-ipd}).
\end{profundiza}

%% ---------------------------------------------------------------------
\section{Para llevar}
\label{sec:re-cierre}
%% ---------------------------------------------------------------------

\begin{parallevar}
\begin{enumerate}
  \item La función de transferencia solo muestra la parte controlable y observable. Estabilizable: los modos no controlables son estables; detectable: los no observables lo son. Nunca se cancelan polos inestables.
  \item Retroalimentación de estado: $u=r-\Km x$ y lazo cerrado $\dot x=(\Am-\Bm\Km)x+\Bm r$. Se pueden poner los $n$ polos donde se quiera si y solo si $\operatorname{rango}\Ctrb=n$, con $\Ctrb=[\Bm\;\Am\Bm\;\cdots\;\Am^{n-1}\Bm]$.
  \item $\Km$ por comparación de coeficientes o por Ackermann: $\Km=[0\cdots0\;1]\,\Ctrb^{-1}q(\Am)$. Las ganancias dependen de la representación; los polos, no.
  \item Polos: dos dominantes por las especificaciones y los demás de cinco a diez veces más lejos (o sobre un cero). Más lejos cuesta más esfuerzo.
  \item Observador: $\dot{\est{x}}=(\Am-\Lm\Cm)\est{x}+\Bm u+\Lm y$; el error cumple $\dot\errest=(\Am-\Lm\Cm)\errest$. Requiere $\operatorname{rango}\Obsv=n$. $\Lm$ se diseña como $\Km$ en el sistema dual.
  \item Separación: los polos del lazo con observador son los de control más los del observador. El compensador resultante es una función de transferencia $U/Y$.
  \item Acción integral: $\dot z=r-y$, $u=-\Km x+K_ez$. Un polo más y error nulo ante escalones.
\end{enumerate}
\end{parallevar}

\begin{profundiza}[Correcciones respecto a la versión anterior del capítulo]
(i) El estimado del observador se escribe $\est{x}$ y el error de estimación $\errest$ (antes $\hat{x}$ y $e$, que chocaban con la variable desviación y con el error de seguimiento). (ii) ``Ackerman'' $\to$ Ackermann. (iii) En el servo de tercer orden, el polinomio es $\lambda^3+(3+K_3)\lambda^2+(2+K_2)\lambda+K_1$ (decía $-\lambda(2-K_2)$). (iv) La ubicación de polos da un $\Km$ \emph{único} solo con una entrada. (v) El ejemplo $\operatorname{diag}(-1,-3)$ ahora tiene solución: $\Km=[0.5\;\;-0.5]$. (vi) $K_e=18\,314$ con $\wn^2$ exacto. (vii) Se agregó la prueba PBH y la noción cuantitativa de controlabilidad. Todos los números se verificaron con \texttt{herramientas/cap09\_numeros.py} y \cmd{verificarCap09.m}.
\end{profundiza}

%% ---------------------------------------------------------------------
\section{Problemas}
\label{sec:re-problemas}
%% ---------------------------------------------------------------------

Los problemas marcados con \etiquetaM{} requieren \matlab/\simulink; los marcados con \etiquetaEstrella{} son de profundización.

\subsection*{Controlabilidad y observabilidad}

\begin{problema}[estructura]\label{prob:re-estructura}
Sean (i) $\Am=\operatorname{diag}(-1,-3)$, $\Bm=[1\;\;1]^\top$, $\Cm=[1\;\;0]$ y (ii) $\Am=\operatorname{diag}(-1,-2,-3,-4)$, $\Bm=[1\;\;0\;\;1\;\;0]^\top$, $\Cm=[1\;\;1\;\;0\;\;0]$. (a) Para (i), construya $\Ctrb$ y $\Obsv$ y concluya. (b) Para (i), halle $\Km$ para polos en $\{-2,-2\}$ y simule desde $x(0)=[1\;\;1]^\top$. (c) El sistema (ii) es la forma modal del ejemplo~\ref{ej:re-friedland}: verifique con los rangos que exactamente dos modos son controlables y dos observables, y diga cuáles. (d) Calcule su función de transferencia y relacione su orden con (c). (e) ¿Es estabilizable? ¿Detectable? (f) Reemplace el valor propio $-2$ por $+2$ y repita (e). ¿Qué consecuencia práctica tiene?
\end{problema}

\begin{problema}[controlabilidad cuantitativa]\label{prob:re-cuantitativa}
\etiquetaEstrella{} Para el sistema del ejemplo~\ref{ej:re-motivador}: (a) calcule $\det\Ctrb$ en función de $a$. (b) Con $a=0$, identifique el modo no controlable escribiendo las ecuaciones y concluya si es estabilizable. (c) Con $a=0$ y $a_{11}=+1$, muestre que \emph{ningún} $\Km$ estabiliza el sistema. (d) Aplique la prueba PBH a (b) y (c). (e) Con polos deseados en $-2$ y $-3$, evalúe $\|\Km\|$ y el número de condición de $\Ctrb$ para $a=2$, $1$, $0.1$ y $0.01$. (f) \etiquetaM{} Simule desde $x(0)=[1\;\;1]^\top$ y reporte el pico de $|u|$. ¿Por qué ``controlable sí o no'' no basta para juzgar un diseño?
\end{problema}

\subsection*{Ubicación de polos}

\begin{problema}[Ackermann y especificaciones]\label{prob:re-ackermann}
(a) Para $1/(s(s+3))$ en variables de fase, verifique controlabilidad y calcule $\Km$ a mano con Ackermann para polos en $-4\pm j4$ (reporte $\Ctrb^{-1}$ y $q(\Am)$); verifique con \cmd{acker} y \cmd{place}. (b) Para $1/(s(s+1)(s+2))$ con la configuración de la \cref{fig:re-servo}, diseñe $\Km$ para $\OS\leq5\%$ y $\ts=2$\,s, con el tercer polo diez veces más a la izquierda. (c) \etiquetaM{} Simule y reporte el pico de $u$. (d) Repita con el tercer polo a $5$ y $50$ veces y haga una tabla con $\OS$, $\ts$ y $|u|_{\max}$. ¿Qué precio se paga por alejarlo?
\end{problema}

\begin{problema}[dos representaciones]\label{prob:re-representaciones}
Para $G_p(s)=\dfrac{10}{(s+10)(s^2+0.5s+1)}$ (ejemplo~\ref{ej:re-trespolos}): (a) \etiquetaM{} reproduzca $\Km$ con \cmd{tf2ss} y diga en qué forma está; (b) repita en variables de fase y explique por qué $\Km$ cambia y los polos no; (c) calcule el error ante escalón y agregue el integrador con un cuarto polo en $-50$.
\end{problema}

\subsection*{Observadores y separación}

\begin{problema}[observador]\label{prob:re-observador}
Para $\Am=\begin{bmatrix}0&1\\-2&-3\end{bmatrix}$, $\Cm=[1\;\;0]$: (a) verifique observabilidad; (b) halle $\Lm$ para polos dobles en $-10$, por coeficientes y por Ackermann; (c) \etiquetaM{} simule planta y observador desde $x(0)=[1\;\;0]^\top$, $\est{x}(0)=0$, sin entrada, y mida cuándo el error cae por debajo del $2\%$; (d) repita con polos en $-4$, $-25$ y $-50$; (e) \etiquetaM{} agregue a $y$ ruido de desviación $0.01$ y mida la desviación del error de estimación en régimen permanente para cada observador. Grafíquela contra la ubicación de los polos y discuta.
\end{problema}

\begin{problema}[compensador completo]\label{prob:re-compensador}
Con la planta del problema anterior, diseñe $\Km$ para polos en $-3\pm j3$ y un observador con polos dobles en $-15$. (a) Halle $\Km$ y $\Lm$. (b) Arme la matriz \eqref{eq:re-aumentado} y verifique que sus valores propios son la unión de los dos conjuntos. (c) Obtenga $U(s)/Y(s)$ con \eqref{eq:re-compensador}: ¿a qué compensador clásico se parece? (d) \etiquetaM{} Simule con $u=-\Km x$ y con el observador partiendo de $\errest(0)=[1\;\;-1]^\top$, y cuantifique el deterioro. (e) Verifique que $\errest(t)$ no depende de la referencia.
\end{problema}

\begin{problema}[robustez del compensador]\label{prob:re-robustez}
\etiquetaEstrella{}\etiquetaM{} Para la planta del problema~\ref{prob:re-representaciones}, con el integrador y un observador cinco veces más rápido que los polos de control: (a) verifique la separación; (b) suponga que el polo real verdadero está en $-8$ y no en $-10$, pero $\Km$ y $\Lm$ se calcularon con el modelo nominal: ¿sigue estable? Encuentre el rango de ubicaciones de ese polo con el que lo es. (c) Perturbe ahora la ganancia de la planta y halle el margen de ganancia efectivo, con y sin observador. ¿Qué dice esto sobre el principio de separación?
\end{problema}

\subsection*{Acción integral}

\begin{problema}[integrador con un cero]\label{prob:re-integralcero}
Para $\Am=\begin{bmatrix}0&1\\-7&-9\end{bmatrix}$, $\Bm=[0\;\;1]^\top$, $\Cm=[4\;\;1]$: (a) verifique controlabilidad y observabilidad y halle $Y/U$ y su cero. (b) Diseñe $u=r-\Km x$ para $\OS=10\%$ y $\tp=2$\,s y calcule el error ante escalón. (c) Agregue el integrador y escriba el polinomio aumentado: ¿por qué $K_e$ aparece en dos coeficientes? (d) Ubique el tercer polo sobre el cero y halle $\Km$ y $K_e$. (e) \etiquetaM{} Verifique que el error es nulo y compare con (b). (f) Repita (d) con el tercer polo en $-8$ y en $-20$ y explique por qué cancelar el cero da el mejor transitorio.
\end{problema}

\begin{problema}[la conexión con el PID]\label{prob:re-ipd}
Para $G_p=\dfrac{1}{(s+1)(s+2)}$ en variables de fase ($x_1=y$, $x_2=\dot y$): (a) escriba $\Am$, $\Bm$, $\Cm$; (b) arme el sistema aumentado con el integrador y su polinomio en función de $K_1$, $K_2$, $K_e$; (c) diseñe para $\OS=10\%$, $\ts=2$\,s y un tercer polo cinco veces más lejos; (d) escriba la ley en el tiempo con $x_1=y$, $x_2=\dot y$, $z=\int e$ y muestre que es un I-PD, identificando $K_c$, $K_i$ y $K_d$; (e) \etiquetaM{} compruébelo numéricamente; (f) ¿qué aporta la formulación de estado si la planta es de orden mayor o si $\dot y$ no se mide?
\end{problema}

%% ---------------------------------------------------------------------
\section{Soluciones seleccionadas}
%% ---------------------------------------------------------------------

\paragraph{Problema~\ref{prob:re-estructura}}
(a) $\Ctrb=\begin{bmatrix}1&-1\\1&-3\end{bmatrix}$ (rango 2: controlable); $\Obsv=\begin{bmatrix}1&0\\-1&0\end{bmatrix}$ (rango 1: el modo $-3$ no se ve). (b) $\Km=[0.5\;\;-0.5]$. (c) Controlables: $-1$ y $-3$; observables: $-1$ y $-2$. (d) $1/(s+1)$: solo el modo controlable y observable. (e) Estabilizable y detectable (los modos escondidos son estables). (f) El modo $+2$ es observable pero no controlable: el sistema es detectable pero \emph{no estabilizable}; ningún controlador lo estabiliza, y la salida crecerá aunque $u$ haga lo que haga.

\paragraph{Problema~\ref{prob:re-cuantitativa}}
(a) $\det\Ctrb=-a$. (b) Con $a=0$, $\dot x_1=-x_1$: el modo $-1$ no es controlable, pero es estable: estabilizable. (c) Con $a_{11}=+1$ el polinomio es $(\lambda-1)(\lambda+2+K_2)$ para todo $\Km$: no estabilizable. (e) $\|\Km\|=4.47$, $5.39$, $23.1$ y $203$; número de condición $4.3$, $5.8$, $50$ y $500$. (f) Pico de $|u|$: $6$, $7$, $25$ y $205$.

\paragraph{Problema~\ref{prob:re-ackermann}}
(a) $\Km=[32\;\;5]$ (pausa de la \cref{sec:re-ackermann}). (b) $\zeta=0.690$, $\wn=2.90$, polos $-2\pm j2.10$ y $-20$: $\Km=[168.0\;\;86.4\;\;21]$. (d) Con el tercer polo a $5$, $10$ y $50$ veces: $\OS=4.73$, $4.94$ y $5.00\%$; $\ts=2.17$, $2.12$ y $2.08$\,s; $|u|_{\max}=84$, $168$ y $840$. Alejar el polo casi no mejora la respuesta y multiplica el esfuerzo.

\paragraph{Problema~\ref{prob:re-observador}}
(b) $\Lm=[17\;\;47]^\top$. (d) Con polos dobles en $-4$, $-25$ y $-50$: $\Lm=[5\;\;-1]^\top$, $[47\;\;482]^\top$ y $[97\;\;2207]^\top$: la ganancia crece con el cuadrado de la velocidad, y con ella la amplificación del ruido.

\paragraph{Problema~\ref{prob:re-compensador}}
(a) $\Km=[16\;\;3]$ y $\Lm=[27\;\;142]^\top$. (b) $\{-3\pm j3,\,-15,\,-15\}$.

\paragraph{Problema~\ref{prob:re-integralcero}}
(a) $\dfrac{s+4}{s^2+9s+7}$, cero en $-4$. (b) $\zeta=0.591$, $\wn=1.95$: $\Km=[-3.21\;\;-6.70]$ y valor final $1.055$ (error $-0.055$). (c) $\lambda^3+(9+K_2)\lambda^2+(7+K_1+K_e)\lambda+4K_e$: $K_e$ entra por los dos elementos de $\Cm$. (d) $\Km=[2.21\;\;-2.70]$, $K_e=3.79$; $\OS=10\%$ y $\tp=2$\,s exactos. (f) Con $-8$: $\OS=11.4\%$, $\tp=1.82$\,s; con $-20$: $11.7\%$ y $1.73$\,s. Cancelando el cero, el lazo es exactamente el segundo orden especificado.

\paragraph{Problema~\ref{prob:re-ipd}}
(b) $\lambda^3+(3+K_2)\lambda^2+(2+K_1)\lambda+K_e$. (c) $\zeta\wn=2$, $\wn=3.38$; con el polo en $-10$: $K_1=49.45$, $K_2=11$, $K_e=114.46$. (d) $u=K_e\int_0^t e\,\dif\tau-K_1y-K_2\dot y$: I-PD con $K_i=K_e$, $K_c=K_1$ y $K_d=K_2$ (proporcional y derivativa sobre la salida).

\respuestaspausas

